\subsection{我们该如何打破恶性循环？}

如果你想打破恶性循环，不要先和这个循环争辩。
先去改变那些会把它一遍遍重新造出来的条件。
在这套框架里，恶性循环之所以持续，是因为当前行为框架总能在关键时刻击败目标状态：
要么障碍太高，要么决策窗口被错过，要么意志能量花在了错误动作上。
实践上的答案，是降低进入下一目标状态的障碍，在真实可用的决策窗口里行动，并安装一个小改变，让下一次重复变得更便宜。

\begin{enumerate}[nosep]
  \item 用具体语言命名这个循环。反复发生的是什么，何时发生，总是哪个状态赢？
  \item 为下一轮指定目标状态。不要瞄准“彻底变好的人生”，而要瞄准下一个真正能稳住的目标状态。
  \item 找出主导障碍。问题是第一步不清楚、意志能量耗尽、环境敌对，还是拖到了决策窗口已经关闭之后？
  \item 在下一次尝试前先降低一个障碍变量。移动物件、预写第一句、换房间，或把进入动作缩小，让目标状态更容易进入。
  \item 在决策窗口里用一个可见动作立刻开始。重点不是一次性解决整个循环，而是阻止旧框架重新把自己搭起来。
  \item 重复同一套低摩擦设置，直到替代模式开始拥有自己的历史。恶性循环会在“新状态比旧状态更容易启动”时开始变弱。
\end{enumerate}

\begin{remark}
一个总是推迟学习的学生，可能会把问题理解为道德意志不够强。
从机制上看，更有用的诊断通常是：目标状态过于模糊，障碍太高，决策窗口被拿去想来想去，而不是被拿去进入。
修复往往很小：低谷来临前先打开材料，用一句话写清第一个任务，并在窗口还没关闭时就开始。
\end{remark}